
\documentclass[11pt]{article}
\usepackage{amsmath,amssymb,amsthm,mathtools}
\usepackage[margin=1in]{geometry}

\newtheorem{theorem}{Theorem}
\newtheorem{lemma}{Lemma}
\newtheorem{proposition}{Proposition}
\newtheorem{definition}{Definition}

\title{Rank-Two Affine S-Unit Prime Production and Erdős--Graham 203}
\author{Jared Wilder}
\date{\today}

\begin{document}
\maketitle

\begin{abstract}
We reduce Erdős--Graham 203 to a finite analytic package over the rank-two
S-unit orbit generated by \(2\) and \(3\).  The elementary algebra, collision
identity, difference-vector multiplicity formula, and final parity-breaking
sieve endpoint are made explicit.
\end{abstract}

\section{Statement}

\begin{theorem}[EG203]
For every \(m\ge 1\) with \((m,6)=1\), there exist \(k,\ell\ge 0\) such that
\[
m2^k3^\ell+1
\]
is prime.
\end{theorem}

\section{Finite reductions}

\begin{lemma}[Type-I congruence]
If \(p\mid mA+1\), then \(mA\equiv -1\pmod p\).  If \(p\) is prime, then
\((p,m)=1\).
\end{lemma}

\begin{lemma}[Type-II S-unit difference]
If
\[
d\mid m2^k3^\ell+1,\qquad d\mid m2^{k'}3^{\ell'}+1,\qquad (d,m)=1,
\]
then
\[
d\mid 2^k3^\ell-2^{k'}3^{\ell'}.
\]
\end{lemma}

\section{Collision identity}

Let \(T_D=\{(k,\ell):k+\ell\le D\}\).  Define
\[
C(D,q)=\#\{(x,y)\in T_D^2:2^{x_1}3^{x_2}\equiv2^{y_1}3^{y_2}\pmod q\}.
\]
For \(h=(a,b)\ne(0,0)\), put
\[
R_{a,b}=\left|2^{a_+}3^{b_+}-2^{a_-}3^{b_-}\right|.
\]
Let
\[
N_D(a,b)=\#\{(x,y)\in T_D^2:x-y=(a,b)\}.
\]
Then
\[
C(D,q)-|T_D|=\sum_{(a,b)\ne(0,0)}N_D(a,b)\mathbf 1_{q\mid R_{a,b}}.
\]

Moreover, if
\[
S=\min(D,D-a-b)-\max(0,-a)-\max(0,-b),
\]
then
\[
N_D(a,b)=
\begin{cases}
0,&S<0,\\
(S+1)(S+2)/2,&S\ge0.
\end{cases}
\]

\section{Analytic package}

The proof is reduced to the following six estimates.

\subsection{Vaughan--Heath-Brown decomposition}
Apply a prime-detecting identity to
\[
\sum_{(k,\ell)\in T_D}\Lambda(m2^k3^\ell+1).
\]

\subsection{Type I}
Use
\[
m2^k3^\ell\equiv -1\pmod p
\]
to control linear congruence sums over the orbit.

\subsection{Type II}
By the pair-GCD lemma, off-diagonal common divisors reduce to divisors of
\(R_{a,b}\).  The Type-II task is to prove a power saving for the weighted
divisor sum
\[
\sum_{(a,b)\ne(0,0)}N_D(a,b)
\sum_{\substack{q\in \mathcal Q\\q\mid R_{a,b}}}\frac1q.
\]

\subsection{Short-relation lattice}
Short relation vectors satisfying \(q\mid R_{a,b}\) are handled explicitly by
the closed formula for \(N_D(a,b)\).

\subsection{Concentration}
For \(H_p=|\langle2,3\rangle\subset \mathbb F_p^\times|\), prove a uniform
estimate of the shape
\[
\sum_{p\le z}\frac1{H_p}\ll \log\log z+O(1).
\]

\subsection{Brun--Hooley parity breaking}
Apply a lower-bound sieve at level \(r\ge\omega(z)+1\) to obtain
\[
\#\{(k,\ell)\in T_D:m2^k3^\ell+1\text{ prime}\}>0.
\]

\section{Lean endpoint}
The corresponding Lean endpoint is:
\[
\texttt{EG203.AnalyticAPI.FullAnalyticPackage}.
\]
Once this is proved, \(\texttt{EG203Closed}\) follows immediately.
\end{document}
